%% notes-tex/microbe-perturb-seq/sections/3-method-explainers.tex
%% SOURCE: hand-authored prose. Protocol steps transcribed from the mirrored
%% papers' Methods sections; every performance number carries its paper.md line
%% in the surrounding text or in experiments/019-perturb-seq-costing/scripts/method_data.py.

\section{Method Descriptions\secstatus{todo}}
\label{sec:methods}

Two families are live options for us. Split-pool is the one we currently prefer,
because it is instrument-free, its cost does not scale with cells, and it has a
credible path to bacteria. Droplet is the incumbent, is available as a service at
UIUC, and recovers roughly five times more mRNA per cell.

\subsection{Yeast-Optimised SPLiT-seq (Brettner et al., 2024)\secstatus{todo}}
\label{sec:brettner}

Brettner et al.\ \citep{brettnerUltraHighthroughputMassively2024} adapted Split-Pool Ligation-based
Transcriptome sequencing to yeast, profiling 43{,}388 cells of \org{S.\ cerevisiae}
(haploid and diploid) and \org{C.\ albicans}. It is the reference protocol for the
split-pool arm, so it is worth walking element by element --- each step exists to
solve a specific failure, and the reasons are what transfer to a new organism.

\notesec{1. Fix before you digest} Cells are fixed in 4\% formaldehyde for
$\sim$18\,h at 4\,\textdegree C, then quenched into 100\,mM Tris with RNase
inhibitor. Formaldehyde crosslinks protein and nucleic acid, which does two jobs
at once: it freezes the transcriptome, and it makes the cell rigid enough to
survive as a reaction vessel through the later handling. The ordering is not
arbitrary. Cell-wall digestion itself triggers the cell-wall-integrity and HOG
stress pathways and induces stress-responsive genes within
minutes~\citep{nadal-ribellesRiseSinglecellTranscriptomics2024}, so a protocol
that digested first would measure its own sample preparation. Fixation first is
what makes the measurement about the perturbation.

\notesec{2. Open the wall, then the membrane, separately} Zymolyase at
0.005\,U/\textmu L in 1\,M sorbitol with 0.1\,M Na$_2$EDTA, 37\,\textdegree C for
15\,min, digests the $\beta$-glucan layer; the sorbitol is osmotic support, since
a wall-less spheroplast bursts in a normal buffer. Then 0.4\% Triton X-100 on
ice for 3\,min permeabilises the membrane. The two steps are tuned against
opposite failures: too little and reagents cannot get in, too much and the cell
stops retaining its own cDNA and the whole single-cell premise collapses. Brettner et al.\ arrived at these values by screening 16 fixation/permeabilisation combinations in
parallel --- itself an argument for the method, since multiplexing let them test
the parameter space in one run rather than 16.

\notesec{3. Round 1 --- barcoded reverse transcription} Cells are counted, diluted,
and distributed into 48 wells of a 96-well plate at $\sim$5{,}000 cells per well.
Each well contains a well-specific barcoded primer pair: a barcoded 15-dT primer
and a barcoded random hexamer, both at 3\,\textmu M, with Maxima H Minus reverse
transcriptase and 7.5\% PEG-6000 as a crowding agent to drive the reaction at low
concentration. The dual priming is a deliberate trade and the single biggest lever
on cost. Oligo-dT captures polyadenylated mRNA specifically; random hexamers give
better 5$'$ coverage but prime on anything, including rRNA. The consequence is
measured: recovered transcripts are 93.7\% rRNA and only 5.75\% mRNA. For a
3$'$-counting screen the 5$'$ coverage buys us nothing, which makes the
hexamer:dT ratio the first thing to retune (Sec.~\ref{sec:bottlenecks}).

\notesec{4. Rounds 2 and 3 --- split-pool ligation} Cells are pooled, filtered,
and redistributed into a fresh 96-well plate. Each well holds a barcode oligo
pre-annealed to a \term{linker strand}, a splint that positions the barcode
against the 5$'$ end of the existing cDNA so T4 DNA ligase can seal the nick.
After 30\,min at 37\,\textdegree C a \term{blocking strand} is added, which
outcompetes the linker and retires that well's barcode so it cannot participate
in any later round --- without it, leftover oligos would scramble the barcode
sequence. Pool, split, repeat for round 3, then terminate with EDTA. Three rounds
of 96 give $96^3=884{,}736$ distinct paths; the round-3 oligo also carries the
10\,nt UMI and a 5$'$ biotin.

\notesec{5. Sublibraries, lysis, and a selection step} Cells are pooled and split
into sublibraries of 5{,}000--20{,}000, which can be frozen. Lysis is SDS plus
proteinase K at 55\,\textdegree C for 2\,h. The biotin now earns its place:
streptavidin beads capture only molecules that received the round-3 oligo, so
anything that fell out of the barcoding scheme is washed away rather than
sequenced. A template-switch reaction on the beads adds a common handle to the
other end, giving both PCR primer sites, and the library is amplified, fragmented,
adapter-ligated, and indexed. That final index is barcode round 4 in all but name
(Sec.~\ref{sec:collisions}).

\begin{keybox}
\textbf{What it delivers.} 100--600 genes and 120--700 mRNA UMIs per cell; 25--50\%
of input cells recovered; barcode collisions well under 5\%. Reagent cost is about
\$0.005 per cell, but 94\% of reads are rRNA and $\sim$75\% of sequenced cells are
discarded --- so the cost that matters is 20$\times$ higher
(Sec.~\ref{sec:application}).
\end{keybox}

%% TODO(library sweep) -- the chemistry drawn here is Brettner's yeast adaptation
%% cross-read against Gaisser's microSPLiT protocol. Sweep the library for the
%% original SPLiT-seq paper (Rosenberg 2018) and Parse's current commercial
%% chemistry, and reconcile: if the commercial kit has moved the UMI or added a
%% round, this schematic silently goes stale.
%% TODO(quote provenance) -- every element of the assembled strand needs a
%% verbatim quote in read_structure.py naming its length and position. The read-2
%% ORDER is the one that already bit us (see the note below the figure); the
%% linker lengths (30 nt and 22 nt, inferred from the STARsolo gaps rather than
%% stated) are the next thing to source directly.
\tcfigfit[0.92]{figures/splitseq-barcoding.pdf}{%
  \textbf{Cell identity in split-pool barcoding is assembled onto the transcript
  itself, which is what removes the need to ever physically isolate a cell.}
  Seven strand states in the yeast-optimised SPLiT-seq workflow of
  \cref{sec:brettner}~\citep{brettnerUltraHighthroughputMassively2024}, each arrow
  annotated with the chemistry that produces the
  next. Strands are drawn 5$'$ left to 3$'$ right, so the barcode block grows
  leftward with each round while the cDNA body stays fixed --- the visual point
  being that a cell's address is written in three instalments, and a cell is
  identified by the \emph{path} it took through three 96-well plates rather than
  by ever having been alone in a well. That is the whole economic argument for the
  method: $96^{3}$ distinct paths cost three plates, not $10^{6}$ containers.
  Two states are drawn double-stranded because that is where the mechanism is
  easy to get wrong. At round 1 the mRNA template lies antiparallel above the
  tracked strand, showing BC1 as an unpaired 5$'$ overhang, and showing that the
  random hexamer rather than the poly-dT primer is the arm that carries the method
  to transcripts without a poly(A) tail --- the single change that lets the same
  chemistry work in
  bacteria~\citep{kuchinaMicrobialSinglecellRNA2021,gaisserHighthroughputSinglecellTranscriptomics2024}.
  At round 2 the linker strand base-pairs underneath
  and a tick marks the nick T4 ligase seals; ligation is templated by that splint,
  which is why an unretired barcode from a previous well would scramble the
  address and why a blocking strand follows every round. The 5$'$ biotin added in
  round 3 is not a label but a selection handle: streptavidin capture at the next
  step discards every molecule that fell out of the scheme, so incompletely
  barcoded material is removed before it can consume sequencing.
  Composed from the protocol as described in \cref{sec:brettner}, not traced from
  any published figure; every parameter shown is sourced there. Drawn at true
  Nature print size in
  \file{notes/assets/drawio/splitseq-barcoding.drawio}.}{fig:splitseq}

\notesec{Reading order follows assembly order} Because each round adds its
barcode outboard of the last, the assembled strand runs UMI, BC3, BC2, BC1, cDNA
from the round-3 end inward, and read 2 reads them in exactly that order. The
STARsolo offsets in Gaisser et al.'s protocol pin it precisely --- UMI at
0--9, BC3 at 10--17, BC2 at 48--55, BC1 at 78--85, linkers filling the gaps ---
which is worth stating because the reverse order is the intuitive guess and it is
impossible: BC1 arrives on the round-1 RT primer, so it can never sit further from
the cDNA than a barcode added later (\cref{tab:read-structure}).

\subsection{Droplet Capture: 10x Chromium with In-Droplet Spheroplasting\secstatus{todo}}
\label{sec:jariani}

Jariani et al.\ \citep{jarianiNewProtocolSinglecell2020} made the 10x Chromium platform work for
yeast with a one-line change, and the reasoning behind it is more interesting than
the change itself.

The obstacle is that the standard 10x workflow lyses cells inside the droplet on a
thermal schedule tuned for a membrane, and a yeast cell wall does not cooperate.
The obvious fix --- spheroplast the cells first, then load them --- is what
Jackson and colleagues did, but it means feeding fragile wall-less cells through a
microfluidic chip, where premature lysis releases mRNA into the shared buffer and
contaminates every droplet downstream. Jariani et al.\ instead put zymolyase \emph{into the
reverse-transcription master mix}, replacing 1\,\textmu L of water with
100$\times$ enzyme, so cells enter the chip walled and intact and are digested only
after they are already isolated. They verified the timing directly: cell counts
hold on ice and through the 6\,min of droplet generation, then drop 200-fold within
20\,min at 53\,\textdegree C. Wall digestion, lysis, and reverse transcription all
happen in the same droplet, in that order, without the cell ever being vulnerable
outside one.

The result was 6{,}118 cells across five conditions with median transcripts per
cell of 528--1{,}261 --- roughly triple Brettner et al.'s per-cell recovery, because
oligo-dT-only capture on a bead is far more mRNA-selective than dT-plus-hexamer in
situ. Library preparation cost fell to \$0.30 per cell against \$4.15 for the
plate-based methods it replaced. Sequencing saturation was only 70--89\%, meaning
deeper sequencing would still have recovered more.

\notesec{The chemistry has moved since} That work used 3$'$ v2. Version 3 roughly
doubles UMI yield on 10x's own matched datasets, and the current generation at the
UIUC core is GEM-X v3.1 on a Chromium X, quoted at $\sim$8\% doublets at 20{,}000
cells recovered with CRISPR guide capture available as a catalogue add-on. Any
comparison drawn from the 2020 paper therefore understates droplet performance
today, and Sec.~\ref{sec:application} uses the newer figures.

%% TODO(quote provenance) -- the 10x row in read_structure.py is the only one with
%% no citation_key/quote: 16 nt CB + 12 nt UMI + dual 10 nt is a vendor spec, not a
%% paper claim. Either mirror the 10x GEM-X user guide and pin a page, or mark the
%% numbers \external{} here. Panel c inherits whatever we decide.
\tcfigfit[0.92]{figures/droplet-10x-barcoding.pdf}{%
  \textbf{Droplet methods buy cell identity with a physical container, which is what
  lets the whole barcode arrive at once --- and what turns the yeast cell wall into a
  scheduling problem rather than a chemistry one.}
  The yeast-adapted 10x Chromium workflow of
  \cref{sec:jariani}~\citep{jarianiNewProtocolSinglecell2020}.
  \textbf{(a)} Cells, beads and oil meet at a junction and partition into droplets.
  Identity here is positional rather than sequential: every molecule tagged inside a
  droplet carries that droplet's bead barcode because the oil wall keeps it that way.
  Loading follows Poisson statistics, so occupancy is kept deliberately low --- most
  droplets are empty, and the doublets that survive that choice are two cells reported
  as one. The container is bought per channel of reagent and its cost scales linearly
  with cells, which is precisely the economics \cref{fig:splitseq} avoids by never
  isolating anything.
  \textbf{(b)} Jariani et al.'s change is one line of the protocol and is entirely
  about \emph{ordering}: 1\,\textmu L of water in the reverse-transcription master mix
  is replaced with 100$\times$ zymolyase, so cells enter the chip walled and are
  digested only once they are already alone. The timing was verified rather than
  assumed --- counts hold on ice and through the 6\,min of droplet generation, then
  fall 200-fold within 20\,min at 53\,\textdegree C, by which point every cell is
  sealed in its own container. The alternative of spheroplasting before loading feeds
  wall-less cells through a microfluidic chip, and mRNA leaking from the ones that
  shear contaminates every droplet downstream, which is a whole-run failure rather
  than a per-cell one.
  \textbf{(c)} The resulting molecule, drawn in the notation of \cref{fig:splitseq}
  so the two can be read against each other. The bead oligo \emph{is} the
  reverse-transcription primer, so as in split-pool the barcode arrives on the primer
  and only the poly-dT tract anneals --- but the entire address, a 16\,nt bead barcode
  plus a 12\,nt UMI, is delivered in one piece instead of over three ligation rounds.
  Once the emulsion is broken the chemistry is bulk, and the read layout that results
  needs 118 cycles against split-pool's 160
  (\file{experiments/019-perturb-seq-costing/scripts/read_structure.py}). That layout
  is the current GEM-X/v3.1 chemistry priced in \cref{sec:application}; Jariani et
  al.\ demonstrated on 3$'$ v2, whose Read 1 is 26\,nt.
  Composed from the protocol as described in \cref{sec:jariani}, not traced from any
  published figure. Drawn at true Nature print size in
  \file{notes/assets/drawio/droplet-10x-barcoding.drawio}.}{fig:droplet}

\subsection{Compressed Perturb-seq (Yao et al., 2024)\secstatus{todo}}
\label{sec:yao}

Yao et al.\ \citep{yaoScalableGeneticScreening2024} is the direct precedent for deliberately
putting \emph{many} perturbations in one cell and recovering individual effects
computationally, rather than treating multi-guide cells as a defect to be filtered
out. It inverts the assumption underneath Sec.~\ref{sec:multiplex}, so it is worth
stating precisely what is compressed and what is assumed.

\notesec{What is compressed} Not the library --- the \emph{samples}. Their framing:
``if the perturbation effects are sparse \dots\ then we can measure a small number
of random composite samples (comprising `linear combinations' of individual sample
phenotypes) and decompress those measurements to infer the effects of individual
perturbations.'' Two routes generate composites: \term{cell-pooling}, deliberately
overloading droplets with several singly-perturbed cells, and \term{guide-pooling},
infecting at high MOI so each cell carries a random set of guides. Guide-pooling is
the one relevant to us. They ran it at MOI 10 in human THP1 cells over 598 target
genes, achieving 2.50 guides per cell, and analysed the 8{,}448 cells carrying
three or more.

\notesec{The decoding} \term{FR-Perturb} factorises the expression matrix by sparse
PCA, runs LASSO on the resulting matrix of perturbations by factors, then multiplies
back out to per-gene effects. The sparsity that makes this work is structural
rather than incidental: regulatory effects are modular, so the effect matrix has
low rank $r$ and few non-zeros $q$ per module, and the required number of composite
samples is $O\!\left((q+r)\log n\right)$ rather than $O(n)$.

\notesec{The assumption, exactly} Guide-pooling ``requires the non-trivial
assumption that the effect sizes of guides tend to combine additively in log
expression space when multiple guides are present in the same cell.'' Interactions
are not modelled in the first-order fit; they are argued to \emph{cancel}, because
each cell gets a random combination, so individual interaction effects average out
``so long as there are an equal number of positive and negative interactions'' ---
a condition they justify by citing genome-wide second-order fitness interactions in
yeast. That is our own data domain, which makes the assumption more checkable for
us than for them.

\notesec{What it delivers, and what it does not} A stated 10-fold cost reduction
for guide-pooling, with performance \emph{increasing} in guides per cell and the
best results in cells carrying four or more --- they found no ceiling within the
range tested and expect further gains from ``even higher efficiency with a greater
degree of overloading.'' The honest limits are equally clear. Pairwise interactions
were not recovered at all: ``none of these effects was significant in either screen
due to insufficient power \dots\ even with a lax significance threshold
$(q<0.5)$.'' And higher-order additivity is essentially untested --- ``very little
is known about the structure of higher-order interaction effects and how they
systematically combine within a cell.''

\begin{keybox}
\textbf{Why this matters here.} It converts multiplexing from a defect into a
design variable, and supplies the constant that Sec.~\ref{sec:multiplex} needs:
$\sim$400 cells per pairwise interaction, four times the 100-cell first-order
floor. But it recovers \emph{main effects} efficiently by assuming interactions
cancel --- and interactions are what metabolic engineering needs. The two goals
pull in opposite directions, which is the tension Sec.~\ref{sec:multiplex}
resolves.
\end{keybox}

\subsection{Ribosomal RNA Depletion\secstatus{todo}}
\label{sec:rrna}

Roughly 85\% of a yeast cell's RNA is ribosomal, and Brettner et al.\ recover
93.7\% rRNA against 5.75\% mRNA. Ninety-four of every hundred reads therefore
buy ribosomal sequence. Since sequencing is 81\% of the split-pool budget
(Sec.~\ref{sec:budget}), this single fraction is the largest cost lever in the
document --- worth about \$99{,}000 at genome scale --- and it is the reason
depletion is treated here as a method component rather than an optimisation.

\notesec{Where it happens in the workflow} This is worth stating precisely,
because there are three plausible points and only one is right. Depletion is
\emph{not} performed on the cell before capture, and \emph{not} after Illumina
indexing. It acts on the \textbf{amplified cDNA library}, after reverse
transcription and the in-cell barcoding rounds, and before fragmentation and
index PCR:

\begin{center}\footnotesize
capture $\rightarrow$ barcode $\rightarrow$ lyse $\rightarrow$ cDNA
$\rightarrow$ \textbf{deplete} $\rightarrow$ fragment $\rightarrow$ index
$\rightarrow$ sequence
\end{center}

Two constraints force that position. Ribosomal RNA cannot be removed before
capture without also stripping the ribosomes whose footprint the in-situ
chemistry depends on, and the cell barcode must already be attached or depletion
would discard cell identity along with the rRNA. Equally it must precede
indexing, because the entire point is to avoid paying to sequence those
molecules. Brandner et al.\ \citep{brandnerPooledSinglecellCRISPRa2025} do exactly this: the
purified cDNA library is incubated with recombinant Cas9 and a pool of guides
targeting rRNA, cleaving rRNA-derived cDNA so it no longer amplifies. Wang et
al.'s M3-seq \citep{wangSinglecellMassivelyparallelMultiplexed2023} places the
same operation at the same point and builds a whole platform around it, which is
the strongest evidence that this position in the workflow is the right one rather
than merely a convenient one.

\notesec{Two independent levers, and they compose} The first is free. Brettner
et al.\ prime with both oligo-dT and random hexamers, and it is the hexamers that
recover rRNA; the 5$'$ coverage they buy is worthless for a 3$'$-counting screen,
so shifting the hexamer:dT ratio costs nothing but a titration. The second is the
Cas9 depletion above, which adds reagents but little labour.

\notesec{How much it is worth, and the caveat} Brandner et al.\ took
\org{E.\ coli} from 4.6\% to 60.2\% mRNA, a 13-fold improvement, and median mRNA
transcripts per cell more than doubled from 54 to 115. But the same protocol in
\org{P.\ putida} moved only 2.2\% to 10.4\%. Depletion efficiency clearly does
not transfer between organisms for free, since the guide pool is designed against
one species' rRNA and against the coverage non-uniformity of one protocol. Yeast
should be better than either bacterium --- it starts at $\sim$85\% rRNA rather
than $>$95\%, and its rRNA genes are fewer and better annotated --- but the
achieved figure has to be measured, and it is the single most valuable number the
pilot can return (Sec.~\ref{sec:bottlenecks}).

\begin{keybox}
\textbf{Treat depletion as mandatory, not optional.} Without it the genome-scale
screen costs \$157\,k and sequencing dominates; with it, \$58\,k and the budget
becomes balanced. No other change in this document moves the total by as much.
\end{keybox}

\subsection{Perturbation Identity Capture\secstatus{todo}}
\label{sec:guide-capture}

Both families above assume the guide can be read out of the cell. In yeast, as our
library stands, it cannot.

Guide RNAs are transcribed from RNA polymerase III promoters --- \gene{SNR52}
upstream, \gene{SUP4} terminator --- so they are neither capped nor polyadenylated
and terminate on a T-tract. A 3$'$ poly(A) capture assay, which is what both
Brettner et al.'s oligo-dT priming and 10x's bead oligo are, sees nothing. Our own MAGIC
library uses exactly this architecture and reads out by plasmid amplicon
sequencing, not RNA~\citep{lianMultifunctionalGenomewideCRISPR2019}; our
reprocessing of it is \file{experiments/016-lian-magic-reprocess}.

The mammalian workarounds do not transfer. 10x direct capture requires a modified
sgRNA scaffold carrying a capture sequence; CROP-seq \citep{datlingerPooledCRISPRScreening2017} relies on the guide cassette
being duplicated into a lentiviral 3$'$ LTR, which is a retroviral trick with no
yeast equivalent.

\notesec{The gap has closed for deletions, and that changes the risk} The 2024
review states plainly that no genome-scale system in yeast lets a genetic
perturbation be traced during
scRNA-seq~\citep{nadal-ribellesRiseSinglecellTranscriptomics2024}. That was true
when written and is no longer, because the same group built one the following
year~\citep{nadal-ribellesSinglecellResolvedGenotypephenotype2025}. The
construction is worth following closely, because it is the design this section
was going to propose.

They rebuilt the yeast knockout collection so its genotype tag is
\emph{transcribed}. In the classical YKOC a \gene{TEF1} promoter drives KanMX and
two 20\,bp barcodes sit outside the transcript, which --- in their words ---
``renders the genotype identity invisible to transcriptomic readouts.'' They
replaced the KanMX marker with \gene{URA3} and shortened the heterologous
terminator from 262\,nt to 43\,nt, which brings the Downtag barcode inside the
\gene{URA3} 3$'$UTR, and added a 5\,nt clone barcode after the stop codon. The
tag is now a polyadenylated Pol II transcript, so an oligo-dT assay sees it. That
is exactly the architecture proposed below, executed at genome scale: 4{,}162
barcoded genotypes, of which more than 3{,}500 were recovered across 1{,}061{,}865
profiled cells, 710{,}952 of them passing QC with a genotype assigned.

Two numbers from it matter more than the rest. Genotype identity was assigned for
\textbf{more than 71\% of cells} --- against the 21--25\% Brandner et al.\ reach
in bacteria --- which is the parameter $q$ that governs everything in
\cref{sec:multiplex}, and it is three-fold better than the figure this document
has been budgeting with. And it was not free: the assignment combines the
transcriptome with a \emph{separate} one-step PCR amplicon library off the
\gene{URA3} 3$'$UTR. Targeted enrichment is doing the work in both systems, so it
should be costed as a second library, not as a primer.

\notesec{What is still open for us} Their perturbation is a chromosomal deletion
with the tag built into the deletion cassette. Ours is a plasmid-borne CRISPRi
guide, so the tag has to ride on the plasmid, and the construct does not exist.
The architecture transfers directly, and its bacterial counterpart is also
precedented: Brandner et
al.\ \citep{brandnerPooledSinglecellCRISPRa2025} built a 316\,bp non-coding
transcript carrying a 12\,bp barcode at its 5$'$ end, a poly(A) tail and a
dedicated enrichment primer site, using an inert gene body as filler, and raised
correct assignment from 5.5\% to 21--25\% by adding the enrichment primer at cDNA
amplification. A design choice worth taking: make the barcode \emph{be} the
guide's own 20\,nt spacer, placed in the 3$'$UTR of a Pol II reporter, so the
guide-to-barcode mapping is known by construction --- that removes a linkage
sequencing run and a whole class of bookkeeping error.

\begin{keybox}
\textbf{Still the critical path, but no longer the risky one.} Every cost in
Sec.~\ref{sec:application} assumes guide identity is recoverable, and for a
plasmid-borne CRISPRi guide it is still unbuilt. What changed is that the same
thing has now been done in yeast at genome scale with a published construct and a
71\% assignment rate, so this is a construction and validation task against a
working template rather than an open research question. The number to reproduce
in the pilot is $q$, because \cref{sec:multiplex} shows it enters the multiplex
ceiling as $q^k$.
\end{keybox}
